\newpage
\phantomsection
\label{prf:finite-discontinuities-darboux-integrable}
\LRAProofFor{thm:finite-discontinuities-darboux-integrable}
\begin{remark*}[Return]\hyperref[thm:finite-discontinuities-darboux-integrable]{Return to Theorem}\end{remark*}
\begin{theorem*}[Bounded Functions with Finitely Many Discontinuities Are Darboux-Integrable]
If \(f:[a,b]\to\mathbb R\) is bounded and has only finitely many
discontinuities, then \(f\) is Darboux-integrable.
\end{theorem*}
\begin{proof}[Professional Standard Proof]
\LRAProofBodyStart
Let \(|f|\le M\), and let the discontinuities be \(c_1,\ldots,c_N\). Fix
\(\varepsilon>0\). Choose open intervals \(G_1,\ldots,G_N\), with
\(c_j\in G_j\), whose total length is less than
\[
  \frac{\varepsilon}{4M+1}.
\]
Let \(G=\bigcup_jG_j\), and let
\[
  K=[a,b]\setminus G.
\]
The function \(f\) is continuous at every point of \(K\). Since \(K\) is
compact, there is \(\delta>0\) such that every subinterval of length less than
\(\delta\) that meets \(K\) has oscillation less than
\[
  \frac{\varepsilon}{2(b-a)}.
\]
Choose a partition \(P\) with mesh less than \(\delta\) and containing all
endpoints of the intervals \(G_j\).

Split the slabs of \(P\) into those contained in \(G\) and those meeting
\(K\). On slabs contained in \(G\), the oscillation is at most \(2M\), and the
total width is less than \(\varepsilon/(4M+1)\), so their contribution to
\(U(f,P)-L(f,P)\) is less than \(\varepsilon/2\). On slabs meeting \(K\), the
oscillation is less than \(\varepsilon/(2(b-a))\), and their total width is at
most \(b-a\), so their contribution is less than \(\varepsilon/2\). Hence
\[
  U(f,P)-L(f,P)<\varepsilon.
\]
By the Darboux criterion, \(f\) is Darboux-integrable.
\end{proof}
\begin{proof}[Detailed Learning Proof]
\LRAProofBodyStart
The proof separates the interval into bad neighborhoods and good regions. The
bad neighborhoods contain the finitely many discontinuities, but they are
chosen with tiny total length. Even if \(f\) oscillates as badly as possible
there, boundedness keeps the total area error small. Outside those
neighborhoods, every point is a continuity point, and compactness turns
pointwise continuity into one uniform oscillation scale. A fine partition then
makes all good slabs nearly flat.
\end{proof}
\begin{remark*}[Proof structure]
Cover finitely many discontinuities by short intervals, use compactness for
uniform oscillation control on the complement, then apply the Darboux
criterion.
\end{remark*}
\begin{dependencies}
\begin{itemize}
  \item \hyperref[def:interval-oscillation-integration]{Oscillation on an Interval}.
  \item \hyperref[def:darboux-sums]{Lower and Upper Darboux Sums}.
  \item \hyperref[thm:darboux-criterion]{Darboux Criterion}.
\end{itemize}
\end{dependencies}
\clearpage
